\documentclass[12pt]{article}


\usepackage{fullpage}
\usepackage{mdframed}
\usepackage{colonequals}
\usepackage{algpseudocode}
\usepackage{algorithm}
\usepackage{tcolorbox}
\usepackage[all]{xy}
\usepackage{proof}
\usepackage{mathtools}
\usepackage{bbm}
\usepackage{amssymb}
\usepackage{amsthm}
\usepackage{amsmath}
\usepackage{amsxtra}
\usepackage{enumitem}
\usepackage{euler}



% Theorem environments
\newtheorem{theorem}{Theorem}[section]
\newtheorem{theorem*}{Theorem}
\newtheorem{definition}[theorem]{Definition}
\newtheorem{corollary}{Corollary}[theorem]
\newtheorem{lemma}[theorem]{Lemma}
\newtheorem{proposition}[theorem]{Proposition}
\newtheorem*{remark}{Remark}
\newtheorem*{claim}{Claim}
\newtheorem*{example}{Example}
\newtheorem*{warning}{Warning}

\newtheorem*{exercisehelper}{Exercise.}
\newenvironment{exercise}[1]{%
  \IfBlankTF{#1}
    {\renewcommand{\exercisehelper}{\textbf{Exercise} \unskip}}
    {\renewcommand\exercisehelper{\textbf{Exercise #1}}}%
  \exercisehelper
}{\endexercisehelper}

\theoremstyle{remark}
\newtheorem*{solution}{Solution}
\newcommand{\mathcat}[1]{\textup{\textbf{\textsf{#1}}}} % for defined terms

\newenvironment{problem}[1]
{\begin{tcolorbox}\noindent\textbf{Problem #1}.}
{\vskip 6pt \end{tcolorbox}}

\newenvironment{enumalph}
{\begin{enumerate}\renewcommand{\labelenumi}{\textnormal{(\alph{enumi})}}}
{\end{enumerate}}

\newenvironment{enumroman}
{\begin{enumerate}\renewcommand{\labelenumi}{\textnormal{(\roman{enumi})}}}
{\end{enumerate}}

\let\H\relax
\let\P\relax
\let\Pr\relax
\let\d\relax
\let\limsup\relax
\let\1\relax
\let\div\relax


\DeclarePairedDelimiter\ceil{\lceil}{\rceil}
\DeclarePairedDelimiter\floor{\lfloor}{ \rfloor}
\newcommand{\fr}[2]{\left(\frac{#1}{#2}\right)}
\newcommand{\norm}[1]{\left\vert\left\vert#1\right\vert\right\vert}
\newcommand{\dbar}{\overline{\partial}}
\newcommand{\d}{\partial}
\newcommand{\RP}{\mathbb{RP}}
\newcommand{\CP}{\mathbb{CP}}
\newcommand{\tbigwedge}{{\textstyle{\bigwedge}}}


\newcommand{\bbni}{\bigbreak \noindent}
\newcommand{\todo}[1]{\textcolor{red}{\textbf{TODO:} #1}}
\newcommand{\red}[1]{\textcolor{red}{#1}}


% Blackboard Font
\newcommand{\N}{\mathbb N}
\newcommand{\Z}{\mathbb Z}
\newcommand{\Q}{\mathbb Q}
\newcommand{\R}{\mathbb R}
\newcommand{\C}{\mathbb C}
\newcommand{\F}{\mathbb F}
\newcommand{\G}{\mathbb G}
\newcommand{\H}{\mathbb H}
\newcommand{\bb}{\mathbb}

% Calligraphic font
\newcommand{\CC}{\mathcal{C}}
\newcommand{\DD}{\mathcal{D}}
\newcommand{\EE}{\mathcal{E}}
\newcommand{\FF}{\mathcal{F}}
\newcommand{\GG}{\mathcal{G}}
\newcommand{\HH}{\mathcal{H}}
\newcommand{\LL}{\mathcal{L}}
\newcommand{\OO}{\mathcal{O}}
\newcommand{\II}{\mathcal{I}}


% Linear Algebra
\DeclareMathOperator{\GL}{GL}
\DeclareMathOperator{\PGL}{PGL}
\DeclareMathOperator{\Mat}{Mat}
\DeclareMathOperator{\Id}{Id}
\DeclareMathOperator{\Tr}{Tr}
\DeclareMathOperator{\tr}{tr}
\DeclareMathOperator{\Nm}{Nm}
\DeclareMathOperator{\opspan}{span}
\DeclareMathOperator{\rk}{rk}
\DeclareMathOperator{\sgn}{sgn}
\DeclareMathOperator{\Sym}{Sym}
\DeclareMathOperator{\Alt}{Alt}
\DeclareMathOperator{\codim}{codim}



% Category Theory / Homological Algebra
\DeclareMathOperator{\id}{id}
\DeclareMathOperator{\colim}{colim}
\DeclareMathOperator{\Aut}{Aut}
\DeclareMathOperator{\End}{End}
\DeclareMathOperator{\Hom}{Hom}
\DeclareMathOperator{\Mor}{Mor}
\DeclareMathOperator{\Ob}{Ob}
\DeclareMathOperator{\Ab}{Ab}
\DeclareMathOperator{\Coh}{Coh}
\DeclareMathOperator{\QCoh}{QCoh}
\DeclareMathOperator{\ann}{ann}
\DeclareMathOperator{\img}{img}
\DeclareMathOperator{\coker}{coker}
\DeclareMathOperator{\Ext}{Ext}
\DeclareMathOperator{\Tor}{Tor}


% Algebraic Geometry
\newcommand{\A}{\mathbb A}
\newcommand{\P}{\mathbb P}
\newcommand{\V}{\mathbb V}
\newcommand{\I}{\mathbb I}
\DeclareMathOperator{\div}{div}
\DeclareMathOperator{\Spec}{Spec}
\DeclareMathOperator{\Proj}{Proj}
\DeclareMathOperator{\Frob}{Frob}
\DeclareMathOperator{\Pic}{Pic}
\DeclareMathOperator{\Weil}{Weil}


% Unclassfified
\DeclareMathOperator{\Cinf}{C^{\infty}}
\DeclareMathOperator{\Ell}{Ell}
\DeclareMathOperator{\CL}{\mathcal{CL}}
\DeclareMathOperator{\Log}{Log}
\DeclareMathOperator{\Arg}{Arg}
\DeclareMathOperator{\Res}{Res}
\DeclareMathOperator{\val}{val}

\DeclareMathOperator{\Supp}{Supp}
\DeclareMathOperator{\cl}{cl}
\DeclareMathOperator{\disc}{disc}
\DeclareMathOperator{\M}{M}
\DeclareMathOperator{\opchar}{char}
\DeclareMathOperator{\argmax}{argmax}
\DeclareMathOperator{\argmin}{argmin}
\DeclareMathOperator{\limsup}{limsup}
\DeclareMathOperator{\Ind}{Ind}
\DeclareMathOperator{\Pr}{\mathbf{Pr}}
\DeclareMathOperator{\Exp}{\mathbb{E}}
\DeclareMathOperator{\Var}{\mathbf{Var}}

\setlength{\parindent}{0pt}
\setlength{\parskip}{5pt}
\setlength{\hfuzz}{4pt}




\begin{document}


\title{Scheme Theory}

\author{Sair Shaikh}
\maketitle




\section{Varieties over a field:}
\begin{definition}
    A variety over a field $k$, or a $k$-variety, is a reduced, separated scheme of finite type over $k$. 
\end{definition}

\begin{definition}
    The category of $k$-schemes, denoted by $\mathrm{Sch}_k$ (where $k$ is a field), is: 
    \begin{itemize}
        \item The objects are morphisms of the form $X \to \Spec k$.
        \item The morphisms are the commutative diagrams: 
        \[  \xymatrix{ X \ar[rd] \ar[rr] & & Y \ar[ld] \\ & \Spec k &} \]
    \end{itemize}
The morphisms $X \to \Spec k$ are called ``structure morphisms''. 
\end{definition}


\begin{proposition}
    The structure morphism give rings of functions, stalks, and residue fields the structure of $k$-algebras. 
\end{proposition}

\begin{proof}
    Let $X$ be a $k$-scheme with structure morphism $\pi: X \to \Spec k$. For any open set $U \subseteq X$, we have a morphism of rings $\pi^\#: k = \OO_{\Spec k}(\Spec k) \to \OO_X(U)$, giving $\OO_X(U)$ the structure of a $k$-algebra. In particular, for any point $p \in X$, we have a morphism of local rings $\pi_p^\#: k = \OO_{\Spec k, \mathfrak{m}} \to \OO_{X,p}$, giving the stalk $\OO_{X,p}$ the structure of a $k$-algebra. Finally, composing with the quotient map $\OO_{X,p} \to \kappa(p)$ gives a morphism $k \to \kappa(p)$, giving the residue field $\kappa(p)$ the structure of a $k$-algebra.    
\end{proof}

\begin{definition}
    A scheme $X$ over a field $k$ is \emph{locally of finite type} over $k$ if it can be covered by affine open sets $\Spec B_i$ where each $B_i$ is a finitely generated $k$-algebra. Moreover, it is \emph{of finite type} if it is also quasicompact, i.e. can be covered by a finite number of such affine opens.
\end{definition}

\begin{definition}
    Let $\pi: X \to Y$ be a morphism of schemes. The diagonal morphism is the unique morphism $\Delta_\pi: X \to X \times_Y X$ whose composition with both projection maps $p_1, p_2$ is the identity map.
\end{definition}

\begin{definition}
    Let $f: X \to Y$. We say that the morphism $f$ is separated if the diagonal morphism $\Delta_f$ is a closed immersion. We also say $X$ is separated over $Y$. A scheme $X$ is separated if it is separated over $\Spec \Z$. 
\end{definition}

\begin{proposition}
    Every morphism of affine schemes is separated.
\end{proposition}

\begin{corollary}
    An arbitrary morphism $f: X \to Y$ is separated if and only if the image of the diagomal morphism is a closed subset of $X \times_Y X$.
\end{corollary}

\section{The Frobenius on a Scheme over $\F_q$:}
\begin{definition}
    For $X/K$ a variety, the degree of a closed point $x \in X_{cl}$ is $\deg(x) = [k(x) : K]$.
\end{definition}

\subsection{Extensions and rational points:}
\begin{proposition}
    Let $L/K$ be a algebraic extension and $X$ an scheme over $K$. Then, 
    \[ X(L) \cong \bigsqcup_{x \in X} \Hom_{K-alg}(k(x), L)\]
    If $X$ is a variety (reduced, seperated, finite type), then this reduces to:
    \[ X(L) \cong \bigsqcup_{x \in X_{cl}} \Hom_{K-alg}(k(x), L)\]
\end{proposition}

\begin{proposition}
    If $X$ is a variety over $K$, then for every $x \in X_{cl}$, the residue field $k(x)/K$ is a finite extension. 
\end{proposition}

\begin{proposition}
    Let $X$ be a variety over $K$ and $x \in X_{cl}$ with $k(x) \cong K[x]/(f(x))$. Then, 
    \[ \Hom_{K-alg}(k(x), L) \leftrightarrow  \{ y \in L : f(y) = 0\}\]
\end{proposition}


\begin{proposition}
    Let $X$ be an affine variety over $K$, i.e. $X - \Spec K[x_1, \cdots, x_n]/(f_1, \cdots, f_k)$. Then an $L$-rational point of $K$ is equivalent to a tuple $(\alpha_1, \cdots, \alpha_n) \in L^n$ such that $f_i(\alpha_1, \cdots, \alpha_n) = 0$ for all $1 \leq i \leq k$.
\end{proposition}

\begin{proposition}
    Let $L/K$ be a field extension. 
    \begin{enumerate}
        \item Taking $L$-points on a $K$-scheme defines a functor $\mathrm{Sch}_K \to \mathrm{Sets}$.
        \item If $U_i$ is an open cover of $X$, then $\bigcup_i U_i(L) = X(L)$. 
        \item For any $K$-schemes $X$ and $Y$, we have $(X \times_K Y)(L) = X(L) \times Y(L)$
    \end{enumerate}
\end{proposition}

\subsection{$\F_q$-rational points and the Frobenius:}
\begin{proposition}
    Let $X$ be a variety over $\F_q$. Then, 
    \[ |X(\F_q)| = \sum_{e | r} e \cdot | \{x \in X_{cl} : \deg(x) = e\}| \]
\end{proposition}

\begin{definition}
    The (absolute) Frobenius morphism on a scheme $X$ over $\F_q$ is denoted $\Frob_{X,q}: X \to X$ and given by the identity on the topological space $X$ and by $f \mapsto f^q$ on the $\OO_X$. It is an $\F_q$-morphism.
\end{definition}

\begin{theorem}
    Let $X$ be a variety over $\F_q$ and fix an embedding $\F_q \hookrightarrow \overline{\F_q}$. The natural embedding $X(\F_{q^r}) \hookrightarrow X(\overline{\F_q})$ identifies $X(\F_{q^r})$ with the fixed points of $\Frob_{X,q}^r$ in $X(\overline{\F_q})$.
\end{theorem}

\subsection{Zeta function:}
\begin{proposition}
    For every variety $X$ over $\F_q$, we have: 
    \[ Z(X, t) = \prod_{x \in X_{cl}} (1 - t^{\deg(x)})^{-1}\]
\end{proposition}

\section{Singular Cohomology:}


\section{Weil Cohomology Theories:}


\end{document}